% This file is part of GUFI, which is part of MarFS, which is released
% under the BSD license.
%
%
% Copyright (c) 2017, Los Alamos National Security (LANS), LLC
% All rights reserved.
%
% Redistribution and use in source and binary forms, with or without modification,
% are permitted provided that the following conditions are met:
%
% 1. Redistributions of source code must retain the above copyright notice, this
% list of conditions and the following disclaimer.
%
% 2. Redistributions in binary form must reproduce the above copyright notice,
% this list of conditions and the following disclaimer in the documentation and/or
% other materials provided with the distribution.
%
% 3. Neither the name of the copyright holder nor the names of its contributors
% may be used to endorse or promote products derived from this software without
% specific prior written permission.
%
% THIS SOFTWARE IS PROVIDED BY THE COPYRIGHT HOLDERS AND CONTRIBUTORS "AS IS" AND
% ANY EXPRESS OR IMPLIED WARRANTIES, INCLUDING, BUT NOT LIMITED TO, THE IMPLIED
% WARRANTIES OF MERCHANTABILITY AND FITNESS FOR A PARTICULAR PURPOSE ARE DISCLAIMED.
% IN NO EVENT SHALL THE COPYRIGHT HOLDER OR CONTRIBUTORS BE LIABLE FOR ANY DIRECT,
% INDIRECT, INCIDENTAL, SPECIAL, EXEMPLARY, OR CONSEQUENTIAL DAMAGES (INCLUDING,
% BUT NOT LIMITED TO, PROCUREMENT OF SUBSTITUTE GOODS OR SERVICES; LOSS OF USE,
% DATA, OR PROFITS; OR BUSINESS INTERRUPTION) HOWEVER CAUSED AND ON ANY THEORY OF
% LIABILITY, WHETHER IN CONTRACT, STRICT LIABILITY, OR TORT (INCLUDING NEGLIGENCE
% OR OTHERWISE) ARISING IN ANY WAY OUT OF THE USE OF THIS SOFTWARE, EVEN IF
% ADVISED OF THE POSSIBILITY OF SUCH DAMAGE.
%
%
% From Los Alamos National Security, LLC:
% LA-CC-15-039
%
% Copyright (c) 2017, Los Alamos National Security, LLC All rights reserved.
% Copyright 2017. Los Alamos National Security, LLC. This software was produced
% under U.S. Government contract DE-AC52-06NA25396 for Los Alamos National
% Laboratory (LANL), which is operated by Los Alamos National Security, LLC for
% the U.S. Department of Energy. The U.S. Government has rights to use,
% reproduce, and distribute this software.  NEITHER THE GOVERNMENT NOR LOS
% ALAMOS NATIONAL SECURITY, LLC MAKES ANY WARRANTY, EXPRESS OR IMPLIED, OR
% ASSUMES ANY LIABILITY FOR THE USE OF THIS SOFTWARE.  If software is
% modified to produce derivative works, such modified software should be
% clearly marked, so as not to confuse it with the version available from
% LANL.
%
% THIS SOFTWARE IS PROVIDED BY LOS ALAMOS NATIONAL SECURITY, LLC AND CONTRIBUTORS
% "AS IS" AND ANY EXPRESS OR IMPLIED WARRANTIES, INCLUDING, BUT NOT LIMITED TO,
% THE IMPLIED WARRANTIES OF MERCHANTABILITY AND FITNESS FOR A PARTICULAR PURPOSE
% ARE DISCLAIMED. IN NO EVENT SHALL LOS ALAMOS NATIONAL SECURITY, LLC OR
% CONTRIBUTORS BE LIABLE FOR ANY DIRECT, INDIRECT, INCIDENTAL, SPECIAL,
% EXEMPLARY, OR CONSEQUENTIAL DAMAGES (INCLUDING, BUT NOT LIMITED TO, PROCUREMENT
% OF SUBSTITUTE GOODS OR SERVICES; LOSS OF USE, DATA, OR PROFITS; OR BUSINESS
% INTERRUPTION) HOWEVER CAUSED AND ON ANY THEORY OF LIABILITY, WHETHER IN
% CONTRACT, STRICT LIABILITY, OR TORT (INCLUDING NEGLIGENCE OR OTHERWISE) ARISING
% IN ANY WAY OUT OF THE USE OF THIS SOFTWARE, EVEN IF ADVISED OF THE POSSIBILITY
% OF SUCH DAMAGE.



\subsubsection{Trace Files}
\label{sec:tracefiles}
Traces files, or flat files, are text files containing
serialized \lstat information pulled from source filesystems. Trace
files allow for easy moving of filesystem metadata to other locations
without needing to copy the potentially massive numbers of directories
and database files found in an index tree. Trace files are generated
by \gufidirtrace and ingested by \gufitraceindex (and \gufitracedir).

In a trace file, each directory is represented by a group of lines
called a ``stanza''. A stanza starts with a line of directory metadata
followed by zero or more lines of regular file and symbolic link
metadata. Stanzas can appear in any order within a trace file and can
be placed in any of the per-thread trace files that are generated.
Traces files can be concatenated together if doing so is necessary or
more convenient.

The columns of a trace file are listed below in the order they appear
in a line. In the code, this data is stored in \texttt{struct~work}
and \texttt{struct~entry\_data} (defined in \texttt{include/bf.h}).

\begin{longtable}{|l|l|p{0.5\linewidth}|}
  \hline
  Name & C Type & Comments \\
  \hline
  name & \texttt{char *} & 4 hex character big endian length + name of entry from the starting path \\
  \hline
  type & \texttt{char} & d[irectory], f[ile], or [sym]l[ink] \\
  \hline
  inode & \texttt{ino\_t} & Stored into \texttt{struct stat} \\
  \hline
  mode & \texttt{mode\_t} & Stored into \texttt{struct stat} \\
  \hline
  nlink & \texttt{nlink\_t} & Stored into \texttt{struct stat} \\
  \hline
  uid & \texttt{uid\_t} & Stored into \texttt{struct stat} \\
  \hline
  gid & \texttt{gid\_t} & Stored into \texttt{struct stat} \\
  \hline
  size & \texttt{off\_t} & Stored into \texttt{struct stat} \\
  \hline
  block size & \texttt{blksize\_t} & Stored into \texttt{struct stat} \\
  \hline
  blocks & \texttt{blkcnt\_t} & Stored into \texttt{struct stat} \\
  \hline
  atime & \texttt{time\_t} & Stored into \texttt{struct stat} \\
  \hline
  mtime & \texttt{time\_t} & Stored into \texttt{struct stat} \\
  \hline
  ctime & \texttt{time\_t} & Stored into \texttt{struct stat} \\
  \hline
  linkname & \texttt{char[4096]} & 4 hex character big endian length + the linkname \\
  \hline
  xattrs & \texttt{struct xattrs} & 4 hex character big endian length + 4 hex character xattr pair count [+ xattr name + \texttt{\textbackslash x1F} + xattr value + \texttt{\textbackslash x1F}]\ldots \\
  & & \\
  & & The 4 hex character big endian length includes the pair count, so it is always at least 4 \\
  & & \\
  & & \texttt{struct xattrs} is defined in \texttt{include/xattrs.h} \\
  \hline
  crtime & \texttt{time\_t} & Creation Time, if available \\
  \hline
  ossint1 & \texttt{int} & OS Specific Integer \\
  \hline
  ossint2 & \texttt{int} & OS Specific Integer \\
  \hline
  ossint3 & \texttt{int} & OS Specific Integer \\
  \hline
  ossint4 & \texttt{int} & OS Specific Integer \\
  \hline
  osstext1 & \texttt{char[1024]} & OS Specific Text \\
  & & 4 hex character big endian length + the string \\
  \hline
  osstext2 & \texttt{char[1024]} & OS Specific Text \\
  & & 4 hex character big endian length + the string \\
  \hline
  pinode & \texttt{long long int} & Parent inode of the current entry \\
  \hline
\end{longtable}

Each of these columns are separated by a fixed (within a single trace
file) delimiter character. The default delimiter is the ASCII Record
Separator character \texttt{\textbackslash x1E}, but can be changed to
any 8-bit character. Characters that can appear in filesystems should
not be used as the delimiter in order to not confuse \gufitraceindex.
With the 4 hex character big endian lengths\footnote{These do not
count against the size of the stored string e.g. an osstext1 of length
1023 (-1 for the \texttt{NULL} terminator) is valid.} prefixed to the
string fields, this should be less of a concern.
